\documentclass[11pt]{article}
\usepackage[a4paper,margin=25mm]{geometry}
\usepackage{amsmath,amssymb}
\usepackage[hidelinks]{hyperref}
\setlength{\emergencystretch}{3em}
\title{Conjecture 00000001294: three real branches every seven levels}
\author{lizaixi01, with OpenAI Codex assistance}
\date{23 October 2026}
\begin{document}
\maketitle
\section{The actual inverse-tree convention}
The claimed bound is false for the standard finite-depth unfolded
inverse tree rooted at $1$, for both normalizations
\[
 T_0(n)=\begin{cases}n/2&n\text{ even},\\3n+1&n\text{ odd},\end{cases}
 \qquad
 T_1(n)=\begin{cases}n/2&n\text{ even},\\(3n+1)/2&n\text{ odd}.\end{cases}
\]
Applegate--Lagarias~\cite{AL}, \S2, pp.~4--5, defines level $N$
by the actual fiber $\{x\ge1:T_1^N(x)=a\}$ and explicitly permits
cyclic-root labels to recur at different levels. A leaf here is a
frontier node at depth $N$ of this finite-depth tree. It is not a
first-hit tree with root-cycle paths deleted. No conclusion about
that different convention, or the Collatz conjecture, is needed.

Let $L_b(N)=\{x\ge1:T_b^N(x)=1\}$ and
$L_b^*(N)=\{x\in L_b(N):3\nmid x\}$. Lean proves
$x\le2T_b(x)$ and hence $x\le2^N T_b^N(x)$.
Thus $L_b(N)$ is the complete actual fiber, represented as a
filtered finite interval $[0,2^N]$, not an assigned branching count.
The pruned set is an actual subset. Every positive nonmultiple
of $3$ stays a positive nonmultiple of $3$ under either forward
map. Therefore all intermediate labels from any $x\in L_b^*(N)$
also survive the $3$-pruning used in~\cite{AL}, equation (2.1).

\section{Three distinct genuine preimages in seven steps}
For every $n>0$ with $3\nmid n$, write $n=9q+r$, where
$r\in\{1,2,4,5,7,8\}$. In the following table each listed pair
$(x,k)$ satisfies $x$ odd and $3x+1=2^k n$:
\[
\begin{array}{c|cc}
r&(x_1,k_1)&(x_2,k_2)\\\hline
1&(12q+1,2)&(48q+5,4)\\
2&(6q+1,1)&(24q+5,3)\\
4&(12q+5,2)&(192q+85,6)\\
5&(24q+13,3)&(96q+53,5)\\
7&(48q+37,4)&(192q+149,6)\\
8&(6q+5,1)&(96q+85,5)
\end{array}
\]
For $T_0$, take the three children
$128n$, $2^{6-k_1}x_1$, $2^{6-k_2}x_2$.
For $T_1$, take $128n$, $2^{7-k_1}x_1$, $2^{7-k_2}x_2$.
Each is positive, not divisible by $3$, and maps to $n$ after
exactly seven actual iterates. The three values are pairwise
distinct. Lean verifies every residue case arithmetically and
proves the iterate identities from $T_b(2x)=x$.

\newpage
\section{True cardinalities and the full asymptotic contradiction}
Iterate this construction from the actual root $1$. Every word
$w\in\{0,1,2\}^m$ determines a genuine label
$v_b(w)\in L_b^*(7m)$. This map is injective: if two final labels
coincide, apply the actual map $T_b^7$ to recover their parents.
Induction identifies their preceding words, and the three-child
injectivity identifies the remaining symbol. In particular,
children of distinct parents cannot collide. The formal finite
type injection into the actual fiber gives
\[
 3^m\le |L_b^*(7m)|\le |L_b(7m)|\qquad(m\ge0,\ b=0,1).
\]
Repeated labels at different levels of the unfolded tree cause
no problem; injectivity is at each fixed level.

Put $c=\sqrt{4/3}$. The exact rational inequality
\[
 (c^7)^2=(4/3)^7<9
\]
gives $0<c^7<3$. Consequently $r=c^7/3$ lies in $(0,1)$ and
$C r^m\longrightarrow0$ for every real $C$.
For every constant $C$ and every cutoff $N$, choose $m\ge N$
with $C r^m<1$. Then $n=7m\ge N$ satisfies
\[
 Cc^n<3^m\le |L_b^*(n)|\le |L_b(n)|.
\]
This refutes every eventual big-$O$ bound, rather than merely
checking finitely many levels or estimating a growth rate from data.

The final Lean statements prove the actual negations
\[
 |L_b^*(n)|\not=O(c^n),\qquad |L_b(n)|\not=O(c^n)
 \quad(b=0,1)
\]
using Mathlib's \texttt{Asymptotics.IsBigO} on the real-valued
cardinality sequences and the actual filter \texttt{atTop}.
The counterexample is for the standard unfolded finite-depth
inverse tree; no claim that every positive integer eventually
reaches $1$ is assumed or proved.

\section{Verification and disclosure}
Lean 4.33.0 and Mathlib commit
\begin{center}\small\texttt{db584cd6d46c92f209a44c0f1c829460d327499d}\end{center}
are pinned. The package supplies a new root build, transitive
axiom audits of every theorem and a full official
\texttt{leanchecker --fresh Main} replay. There is no
\texttt{sorry}, \texttt{native\_decide} or additional axiom.
OpenAI Codex assisted; a separate supervising Codex agent
reviewed the proof, source convention, report and PDF. This is AI
review, not human certification. Organizer acceptance is pending.

\begin{thebibliography}{1}
\bibitem{AL} D.~Applegate and J.~C.~Lagarias,
\emph{Density Bounds for the $3x+1$ Problem I. Tree-Search Method},
December 15, 1993, \S2, pp.~4--5.
\url{https://websites.umich.edu/~lagarias/doc/applegateI.pdf}.
The shortcut map is equation (1.1); the unfolded cyclic-root and
$3$-pruned conventions are explicitly given in \S2.
\end{thebibliography}
\end{document}
